\input{preamble}

\title{Section 4.5 --- Exponential Models}

\begin{document}
\maketitle
\noindent In these exercises, when no rounding instructions are given, round in the way that makes the most sense for the context of the problem.
\section*{Exponential Applications}
\begin{senum}
\item ImplausiCorp's annual profits are growing exponentially.  Each year, they increase by 20\%.  This year, they made \$17 billion dollars.
\begin{clist}
\item Do fractional numbers of years make sense in this problem?
\item Write a function $P(t)$ giving the profits $t$ years from now.
\item What will their profit be 4 years from now?  Round your answer to the nearest 0.01 billion.
\item How many years from now will their profit first exceed \$1 trillion (that is, \$1,000 billion)?
\end{clist}
%%%%%%%%%%%%%%%%%%%%%%%%%%%%%%%%%%%%%%%%%
\item Suppose 10,000 copies of a certain limited edition card are printed for a collectable card game.  Each year, 12\% of the cards in circulation are destroyed.
\begin{clist}
\item Do fractional numbers of years make sense in this problem?
\item Write a function $n(t)$ giving the number of cards that remain $t$ years from now.
\item How many cards will be in circulation after 5 years?
\item When will only 1 card remain in circulation?
\end{clist}
%%%%%%%%%%%%%%%%%%%%%%%%%%%%%%%%%%%%%%%%%
\item Carbon-11, a radioactive isotope, has a half-life of approximately 20 minutes.  Suppose a sample is obtained which contains 27 grams of carbon-11.
\begin{clist}
\item Do fractional numbers of minutes make sense in this problem?
\item Write a function $m(t)$ giving the mass of Carbon-11 that will remain after $t$ minutes.
\item How much Carbon-11 remains after 1 hour? Round to two decimal places.
\item The research facility determines that the sample will be safe to dispose of when only 0.2 grams remain. After how long will this sample be safe to dispose of? Answer as a whole number of minutes.
\end{clist}
\item A drug is metabolized with the half-life of 7.3 hours. An initial dose of 200mg is initially administered.
\begin{clist}
\item Write a function $m(t)$ giving the amount of the drug remaining in the patient's system $t$ hours after it is administered.
\item In order to remain effective, the next dose must be administered before the amount of the drug in the patient's system drops below 10mg. After how many hours should the next dose be administered? (Answer as a whole number of hours.)
\item At the time you found in part (b), how much of the drug should be administered to bring the amount in the patient's system back up to 200mg? (Answer as a whole number of mg. Your answer should ensure that the amount of the drug in the patient's system does not go \emph{over} 200mg.)
\item If a second dose is administered 24 hours after the first, how much should be given to bring the amount of the drug in the patient's system back up to 200mg? (Answer as a whole number of mg. Your answer should ensure that the amount of the drug in the patient's system does not go \emph{over} 200mg.)
\end{clist}
\clearpage
%%%%%%%%%%%%%%%%%%%%%%%%%%%%%%%%%%%%%%%%%
\item The population of a small town declines exponentially after its primary industry goes out of business. In 1997, the population was 30,850. By 2015, it had fallen to only 9,810. (Hint: in setting up this problem, let $t$ represent the number of years since 1997.)
\begin{clist}
\item Write a function $p(t)$ giving the population $t$ years after 1997.
\item What was the population in 2006?
\item What was the last year that the population was at least 20,000?
\item What percent of the population is lost each year? Answer as a percent rounded to one decimal place.
\end{clist}
\end{senum}
%%%%%%%%%%%%%%%%%%%%%%%%%%%%%%%%%%%%%%%%%
\section*{Compound Interest}
You may wish to reference the compound interest formula in solving these problems:
$$P(t)=P_0\left(1+\frac{r}{n}\right)^{nt}$$
\begin{senum}
\item Julia borrows \$1800 to pay for school at 12\% interest, compounded monthly.  How much will she need to pay back after 4 years?
\item Alan makes an investment in an account that pays 11\% interest, compounded annually.  After 10 years, he wants the account to be worth \$8000.  How much must he invest?
\item Brad borrows \$1500 at 8\% interest, compounded semiannually.  How much must he pay back at the end of 3 years?
\item Carol invests \$740 at 3.5\% interest, compounded quarterly.  How much is the investment worth after 18 \emph{months}?
%%%%%%%%%%%%%%%%%%%%%%%%%%%%%%%%%%%%%%%%%
\item Kyle wins \$17000 on a game show, and decides to save it in an account which pays 3\% interest, compounded quarterly.
\begin{clist}
\item How much will his deposit be worth after 10 years?
\item What is the doubling time on this investment? Round to three decimal places.
\item If he had chosen an account with 4\% interest, compounded annually, by how much would that reduce the doubling time? Round to three decimal places.
\item With each account, how long must he wait to have \$50,000?
\end{clist}
%%%%%%%%%%%%%%%%%%%%%%%%%%%%%%%%%%%%%%%%%
\item A vampire deposits \$1 in a long-term account which pays 8\% interest, compounded annually.
\begin{clist}
\item How much money will he have after 100 years?
\item How long must he wait in order to have \$1,000,000?
\end{clist}
\end{senum}

\end{document}